\documentclass[conference]{IEEEtran}
% IEEEtran ships with TeX Live. If your target venue provides its own
% class file, drop it beside this file and change the line above.

\usepackage{cite}
\usepackage{amsmath, amssymb}
\usepackage{graphicx}
\usepackage{booktabs}

\begin{document}

\title{Your Paper Title}

\author{
  \IEEEauthorblockN{First Author}
  \IEEEauthorblockA{
    \textit{Department} \\
    \textit{Institution} \\
    City, Country \\
    first.author@example.com
  }
  \and
  \IEEEauthorblockN{Second Author}
  \IEEEauthorblockA{
    \textit{Department} \\
    \textit{Institution} \\
    City, Country \\
    second.author@example.com
  }
}

\maketitle

\begin{abstract}
  A single paragraph, typically 150--250 words for IEEE venues: context,
  gap, approach, headline result.
\end{abstract}

\begin{IEEEkeywords}
  keyword one, keyword two, keyword three
\end{IEEEkeywords}

\section{Introduction}
\label{sec:introduction}

IEEE conference papers open by framing the problem and stating the
contribution explicitly. Cite prior art as you go~\cite{knuth1984}.

\section{Related Work}
\label{sec:related}

Group prior work by approach rather than chronology.

\section{Method}
\label{sec:method}

\begin{equation}
  \label{eq:example}
  \hat{y} = \arg\max_{y} P(y \mid x) .
\end{equation}

\section{Experiments}
\label{sec:experiments}

\begin{table}[htbp]
  \caption{Replace with your own results.}
  \label{tab:results}
  \centering
  \begin{tabular}{lrr}
    \toprule
    Method   & Accuracy & Latency (ms) \\
    \midrule
    Baseline & 0.81     & 124          \\
    Ours     & 0.93     & 101          \\
    \bottomrule
  \end{tabular}
\end{table}

\section{Conclusion}
\label{sec:conclusion}

State the contribution and the most promising next step.

\bibliographystyle{IEEEtran}
\bibliography{references}

\end{document}
